% ── Résumé / Abstract ──────────────────────────────────────────────────────────
\newpage
\thispagestyle{empty}
\section*{Résumé}
\addcontentsline{toc}{section}{Résumé}

L'organisation des soutenances en milieu universitaire constitue un processus complexe mobilisant de multiples acteurs -- étudiants, encadrants, membres de jury et personnel administratif -- autour de contraintes temporelles, spatiales et réglementaires. Le présent projet vise à concevoir et développer une application web complète pour la gestion centralisée, automatisée et fiable des soutenances et des jurys au sein de la Faculté des Sciences Ben M'Sick. Le système repose sur une architecture moderne à trois couches : un frontend React, un backend Spring Boot et une base de données MySQL. Il couvre l'ensemble du cycle de vie d'une soutenance, de la planification des créneaux et de l'affectation des jurys jusqu'à la saisie des évaluations et la génération automatique des documents administratifs (convocations, procès-verbaux, fiches d'évaluation). Une attention particulière est portée à la détection des conflits horaires et à l'équilibrage de la charge des enseignants. Ce rapport présente l'étude de l'existant, l'analyse des besoins, la conception UML, l'architecture technique et la modélisation de la base de données du système.

\vspace{0.5cm}
\noindent\textbf{Mots-clés :} Gestion de soutenances, planification, jury, évaluation, Spring Boot, React, JWT, UML, système d'information universitaire.

\vspace{1cm}
\section*{Abstract}
\addcontentsline{toc}{section}{Abstract}

Organizing thesis defenses in a university environment is a complex process involving multiple actors -- students, supervisors, jury members, and administrative staff -- around temporal, spatial and regulatory constraints. This project aims to design and develop a complete web application for the centralized, automated and reliable management of defenses and juries at the Faculty of Sciences Ben M'Sick. The system is based on a modern three-tier architecture: a React frontend, a Spring Boot backend, and a MySQL database. It covers the entire lifecycle of a defense, from slot planning and jury assignment to evaluation entry and automatic generation of administrative documents (summons, minutes, evaluation sheets). Special attention is given to conflict detection and teacher workload balancing. This report presents the study of the current system, requirements analysis, UML design, technical architecture, and database modeling.

\vspace{0.5cm}
\noindent\textbf{Keywords:} Defense management, planning, jury, evaluation, Spring Boot, React, JWT, UML, university information system.
